\section{System Integration and API Gateway}
\label{sec:system_integration}

To build a production-ready application, the isolated reasoning components of Omni-IPS are assembled into a cohesive microservice architecture. The system consists of a high-performance **FastAPI Gateway** backend and an interactive **Streamlit Client** frontend. Communication is managed via strictly typed JSON APIs and reactive asynchronous streams.

\subsection{Omni-IPS Core Technology Stack}
The complete technical ecosystem powering Omni-IPS is selected for speed, rigid typing, and robust database operations. We summarize the core tech stack components in Table~\ref{tab:tech_stack}.

\begin{table}[h]
    \centering
    \caption{Omni-IPS Architectural Technology Stack Summary}
    \label{tab:tech_stack}
    \begin{tabular}{llcl}
        \toprule
        \textbf{Component Layer} & \textbf{Framework/DB} & \textbf{Icon} & \textbf{Functional Role in System} \\
        \midrule
        Core Language & Python 3.11 & \textcolor{blue!70!black}{\faPython} & Logical scanner, parsing, and reasoning chainer core \\
        Backend API & FastAPI & \textcolor{teal!70!black}{\faServer} & Async REST endpoints, Pydantic gates, streaming response \\
        Graph Database & Neo4j DB & \textcolor{blue!60!black}{\faNetworkWired} & Canonical rule \& fact storage, structural schema integrity \\
        Vector Database & Qdrant & \textcolor{red!70!black}{\faBrain} & High-dimensional semantic embeds ($d=384$) \& query matching \\
        Web Frontend UI & Streamlit & \textcolor{orange!80!black}{\faDesktop} & Typewriter markdown rendering, real-time chunk streams \\
        \bottomrule
    \end{tabular}
\end{table}

\subsection{FastAPI API Gateway}
The backend service in \texttt{api/main.py} acts as the orchestrator. It receives client requests, triggers the RAG router, queries Neo4j for active rules, runs the symbolic engine, and returns results. 

The API documentation is automatically generated via OpenAPI (Swagger UI), as shown in Figure~\ref{fig:api_gateway} (sourced from \texttt{assets/omni\_ips\_api\_gateway.png}).

\begin{figure}[h]
    \centering
    \includegraphics[width=0.85\textwidth]{omni_ips_api_gateway.png}
    \caption{FastAPI Interactive Swagger UI API gateway displaying the logical endpoint blocks for the logic solver, GraphRAG inference, and the LLM explainability agent.}
    \label{fig:api_gateway}
\end{figure}

The two primary endpoints exposed by the gateway are:
\begin{enumerate}
    \item \texttt{POST /api/solve}: Synchronously performs the entire Neuro-Symbolic GraphRAG query resolving, returning the structured initial facts, goal facts, rule execution steps, and final working memory state.
    \item \texttt{POST /api/explain/stream}: Asynchronously streams human-readable explanations. It receives the logical trace and utilizes LangChain's \texttt{astream} API to stream response chunks chunk-by-chunk in real-time, preventing high-latency delays.
\end{enumerate}

The implementation of the async explain streaming route is shown in Listing~\ref{lst:explain_stream}.

\begin{lstlisting}[language=Python, style=pythonstyle, label={lst:explain_stream}, caption={FastAPI asynchronous streaming implementation of the explainability agent in \texttt{api/main.py}.}]
@app.post("/api/explain/stream", tags=["Explainability Agent"])
async def explain_proof_stream(request: ExplainRequest):
    """
    Explainability Agent (Streaming): Streams a rich educational 
    explanation chunk-by-chunk using LangChain's async stream API.
    """
    domain = request.domain.lower()
    
    if request.goal_reached:
        prompt_text = PROOF_EXPLANATION_PROMPT.format(...)
    else:
        prompt_text = PROOF_FAIL_PROMPT.format(...)

    async def generate_chunks():
        try:
            llm = get_llm_client()
            chain = LLMChain(llm=llm, prompt=EXPLAIN_PROMPT)
            async for chunk in chain.astream({"prompt_text": prompt_text}):
                yield chunk
        except Exception as e:
            # Fallback to offline template streaming chunks
            for word in OFFLINE_TEMPLATE.split(" "):
                yield word + " "
                await asyncio.sleep(0.02)

    return StreamingResponse(generate_chunks(), media_type="text/plain")
\end{lstlisting}

\subsection{Streamlit Interactive Client}
The frontend interface, written in Streamlit (\texttt{ui/app.py}), is designed for premium user experiences. It uses glassmorphic elements, HSL colors, and smooth rendering states. 
When the user submits a mathematical or chemical query, the UI executes two sequential calls to the backend:
\begin{enumerate}
    \item It issues a post request to \texttt{/api/solve} and renders the initial mapped facts and goal, followed by a chronological execution step trace inside an educational step-by-step expander.
    \item It establishes an HTTP chunked transfer connection to \texttt{/api/explain/stream}. Streamlit captures the yielded text fragments reactively and displays them using a typewriter-like markdown block via \texttt{st.write\_stream()}, providing a natural, lag-free user experience.
\end{enumerate}

Figure~\ref{fig:ui_demo} (sourced from \texttt{assets/demo\_Pythagore\_1.png}) illustrates this interface in action, demonstrating the streaming explanation of the Pythagorean proof.

\begin{figure}[h]
    \centering
    \includegraphics[width=0.85\textwidth]{demo_Pythagore_1.png}
    \caption{Interactive Streamlit web frontend displaying a proof query. The interface displays the mapped facts, the active solver execution steps, and the real-time explanation streaming.}
    \label{fig:ui_demo}
\end{figure}
